For an integer $d\ge 2$ and $k\in\mathbb{N}_0$, define
\[
P(d,k) := \Big\{\sum_{i\ge k} \varepsilon_i d^i : \varepsilon_i\in\{0,1\}\text{ and only finitely many }\varepsilon_i\neq 0\Big\}.
\]
So $P(d,0)$ is the set of nonnegative integers with base-$d$ digits in $\{0,1\}$.

Given integers $3\le d_1<\cdots<d_r$, consider representations of integers of the form
\[
 n = \sum_{j=1}^r c_j a_j, \qquad c_j\in\{0,1\},\ a_j\in P(d_j,k).
\]

**Ambiguity in the statement.** The problem text writes the condition
$\sum_{1\le i\le r} 1/(d_r-1)\ge 1$, which simplifies to $r/(d_r-1)\ge 1$ and conflicts with the later sentence "Pomerance observed that the condition $\sum 1/(d_i-1)\ge 1$ is necessary".

**Minimal correction consistent with the surrounding text:** replace the condition by
\[
\sum_{i=1}^r \frac{1}{d_i-1}\ \ge\ 1.
\]

**Question 1 (k=0).** Under $\sum_{i=1}^r 1/(d_i-1)\ge 1$, must every sufficiently large integer be representable as above with $k=0$?

**Question 2 (k\ge 1).** Under $\gcd(d_1,\dots,d_r)=1$ and $\sum_{i=1}^r 1/(d_i-1)\ge 1$, does there exist $k\ge 1$ such that every sufficiently large integer is representable with this fixed $k$?
